\section{Introduction}
\label{sec:ch4_introduction}

Despite advances in cloud computing~\cite{shang2023online}, optimizing resource utilization in edge-cloud networks remains challenging due to the ever-changing nature of cloud-native workloads~\cite{10646623}. Recently, workload prediction has become crucial for service orchestration and resource management in edge-cloud networks~\cite{10068185}. Accurate workload prediction facilitates the proactive and efficient allocation of computational resources based on real-time requirements, minimizing both over-provisioning and under-provisioning~\cite{10229064, 8258257, 9068614}. However, achieving precise workload predictions is complex due to factors such as changes in user demand and workload migrations~\cite{ouyang2023dynamic}, leading to discrepancies between training models and real-world conditions. This generalization challenge, where models trained on one set of workload traces perform poorly on different traces, represents a critical bottleneck for practical deployments~\cite{zou2024generalization}. These temporally inconsistent events that occur on various timescales degrade prediction accuracy and orchestration efficiency~\cite{10029931}. This issue is particularly challenging for container-level workloads due to poor resource isolation across containers, short container lifespans, high container density, and low data generation per container~\cite{10417087, 10202641}. The growing prevalence of microservice applications, where different microservices are independently deployed across containers requiring separate elasticity operations, further highlights the importance of accurate container-level workload prediction~\cite{10202641, shang2023online, 10052731}.

Given the dynamic nature of cloud-native workloads, efforts have focused on developing machine learning models capable of generalizing well to unseen data. Transfer learning techniques~\cite{wang2019transfer} have shown limited success when source and target distributions differ significantly. Domain adaptation methods work for zero-shot scenarios, but require substantial computation. \acp{RNN} have been explored to predict workloads but face challenges such as decreased memory and limitations in sequential processing~\cite{hochreiter1998vanishing, benidis2022deep}. \acp{CNN} were used to capture local patterns in time-series data, but struggled with long-term dependencies~\cite{acmtimeseriesreview2024}. The recent \ac{ModernTCN}~\cite{ModernTCN} improves upon traditional convolution approaches with larger kernels and separated processing pathways, achieving strong results efficiently, although it lacks specific mechanisms to address the challenges of heterogeneous container workloads. Hybrid models combining \ac{CNN} and \ac{RNN} architectures have been developed to address these issues~\cite{xu2022esdnn}. Ensemble methods~\cite{kim2024ensemble} combine multiple models to improve predictions, but their computational demands limit practicality for resource-constrained edge environments. Integration of \acp{GNN} and \acp{GAN} into \acp{RNN} has shown promise in capturing both spatial and temporal information~\cite{li2024evogwp, esteban2017real, AGCRN}. Despite these advancements, these hybrid models still face limitations inherent to their \ac{RNN} and \ac{CNN} components.

Recent advancements in machine learning for time-series forecasting have introduced techniques that can identify complex dependencies across different timescales~\cite{gort2024forecasting}. Attention mechanisms have demonstrated effectiveness in generalization across diverse time-series patterns~\cite{vaswani2017attention, wen2022transformers}. Attention-enhanced neural networks~\cite{zhang2022attention} and attention-based transfer learning~\cite{li2022robust} have shown particular promise for workload prediction, selectively focusing on critical temporal patterns while filtering out noise. Efficient modeling approaches that minimize computational overhead while maintaining high prediction accuracy are particularly valuable for resource-limited edge-cloud environments~\cite{fasterandlightertransformers, gorbett2023sparse, STRec}. For a detailed review of related literature on cross-application generalization, see Section~\ref{sec:ch2_sota_omni}.

In this chapter, we introduce \ac{OmniFORE}. Our framework addresses four critical generalization challenges that limit existing approaches: (1) computational inefficiency in resource-constrained environments, (2) inability to capture multi-scale temporal patterns, (3) slow adaptation to new workload characteristics, and (4) overfitting to specific workload features. Unlike approaches that merely combine existing techniques, \ac{OmniFORE} introduces a novel temporal clustering method that identifies fundamental workload patterns across heterogeneous traces, enabling effective generalization even to previously unseen workload types. Temporal clustering of time-series enables strategic training on representative subsets from large datasets, capturing both short-term stability and long-term dynamic changes in container-level features. This ensures consistent performance even with previously unseen data. Our method also includes an efficient data sampling strategy that reduces computational overhead, ensuring consistent performance across diverse workload scenarios. Extensive experiments with real-world data validate our approach, demonstrating higher prediction accuracy and faster inference speeds compared to state-of-the-art methods. Our primary contribution is threefold:

\begin{enumerate}
\item We introduce \textbf{a forecasting framework with temporal clustering} to ensure model generalization in edge-cloud networks.
\item We develop \textbf{a data sampling strategy} that minimizes computational overhead while maintaining high performance.
\item We evaluate \ac{OmniFORE} with real-world data, demonstrating \textbf{higher prediction accuracy and faster inference speeds}.
\end{enumerate}

The remainder of this chapter is organized as follows: Section~\ref{sec:ch4_system_model} presents our system model and explains the generalization challenges inherent in edge-cloud networks. Section~\ref{sec:ch4_proposed_solution} presents our novel approach, with particular emphasis on our clustering techniques and the overall framework architecture. Section~\ref{sec:ch4_evaluation} offers an analysis of our experimental results. Finally, Section~\ref{sec:ch4_conclusion} concludes the chapter with our key findings.

